%nt 6
Dag einn setjast $n$ höfðingjar við hringborð og skiptast á að vera konungur ríkis. Á fyrsta degi er einhver höfðingi valinn sem kóngur. Á öðrum degi er höfðinginn sem situr beint til vinstri við kónginn valinn sem næsti kóngur. Á þriðja degi er höfðinginn sem situr tveimur sætum til vinstri við kónginn valinn sem næsti kóngur. Áfram gengur þetta svona, svo almennt verður höfðinginn sem situr $m-1$ sæti til vinstri við síðasta konung, kóngurinn á degi $m$. Fyrir hvaða gildi á $n$ munu allir höfðingjarnir einhverntímann verða kóngar?